%%%%%%%%%%  scan idx 265  |  folio 254  %%%%%%%%%%

\subsection*{3. Biextensiones de esquemas en grupos $(P,Q)$ por
$\underline G_m$: generalidades}

En toda esta sección fijamos un esquema de base $S$ y, como se anunció
en la sección 0, trabajamos en el topos fppf de $S$.

\medskip
\noindent 3.1. Recordemos primero el ejemplo más importante de un
grupo $\Biext^1(P,Q;\underline G_m)$: el caso en que $P$ y $Q$ son
dos esquemas abelianos sobre $S$ y $\Biext^1$ se interpreta como el
grupo de correspondencias divisoriales sobre $(P,Q)$ (VII 2.9.4). Los
acoplamientos $\ell$-ádicos de la sección precedente son entonces los
acoplamientos familiares de la teoría de los esquemas abelianos [1].

\medskip
\noindent 3.2. Más generalmente, supongamos solamente que $Q$ es un
esquema abeliano. Es bien sabido (VII 1.3.8) que
\begin{equation}\tag{3.2.1}
 \uHom(Q,\underline G_m)=0,
\end{equation}
y, por tanto, (1.4) da un isomorfismo canónico
\begin{equation}\tag{3.2.2}
 \Biext^1(P,Q;\underline G_m)\simeq\Hom(P,Q^*),
\end{equation}
donde
\begin{equation}\tag{3.2.3}
 Q^*=\uExt^1(Q,\underline G_m)
\end{equation}
es el esquema abeliano dual de $Q$. (Para la representabilidad de
$Q^*$, debida a M. Raynaud en el caso general, véase [1].)

Cuando $P$ es también un esquema abeliano, se tiene, simétricamente a
(3.2.2), un isomorfismo canónico
\begin{equation}\tag{3.2.4}
 \Biext^1(P,Q;\underline G_m)\simeq\Hom(Q,P^*).
\end{equation}

%%%%%%%%%%  scan idx 266  |  folio 255  %%%%%%%%%%

Así, el dato de una biextensión $\xi$ de $(P,Q)$ por $\underline G_m$
puede interpretarse indistintamente como el dato de un homomorfismo
\[
 u:P\longrightarrow Q^*
\]
o de un homomorfismo
\[
 v:Q\longrightarrow P^*.
\]

Recordemos que un homomorfismo de esquemas abelianos
$w:A\longrightarrow B$ es un isomorfismo si y solo si, para todo
entero $n>0$, el homomorfismo inducido
${}_nw:{}_nA\longrightarrow{}_nB$ es un isomorfismo. En el caso de
$u$ y $v$ anteriores, ${}_nu:{}_nP\longrightarrow{}_nQ^*$ se
identifica con el homomorfismo
${}_nP\longrightarrow D({}_nQ)$ definido por el acoplamiento
\[
 \varphi_\xi^n:{}_nP\times{}_nQ\longrightarrow\underline G_m
\]
(véase 2.3.11), mientras que ${}_nv:{}_nQ\longrightarrow{}_nP^*$ se
identifica con el homomorfismo
${}_nQ\longrightarrow D({}_nP)$ definido por el acoplamiento
\[
 \psi_\xi^n:{}_nP\times{}_nQ\longrightarrow\underline G_m.
\]
Ahora bien, $\varphi_\xi^n=-\psi_\xi^n$ (2.1.10), de modo que, salvo
signo, ${}_nu$ y ${}_nv$ son transpuestos uno de otro. Como ${}_nP$ y
${}_nQ$ son finitos y localmente libres sobre $S$, y por tanto
satisfacen la dualidad de Cartier, se sigue que ${}_nu$ es un
isomorfismo si y solo si ${}_nv$ lo es. En consecuencia, $u$ es un
isomorfismo si y solo si $v$ lo es.

Una formulación equivalente, aunque algo menos simétrica, es la
siguiente. Considérese la biextensión canónica (la biextensión de
Weil) de $(P,P^*)$ por $\underline G_m$, correspondiente al
homomorfismo identidad $v:P^*\longrightarrow P^*$. Entonces el
homomorfismo canónico correspondiente
\[
 u:P\longrightarrow P^{**}
\]
es un \underline{isomorfismo}. Este es el \underline{teorema de
bidualidad} de Barsotti y Cartier.
